
\documentclass{article}
\usepackage{colortbl}
\usepackage{makecell}
\usepackage{multirow}
\usepackage{supertabular}

\begin{document}

\newcounter{utterance}

\centering \large Interaction Transcript for game `hot\_air\_balloon', experiment `air\_balloon\_survival\_en\_reasoning off\_hard', episode 5 with gpt{-}oss{-}120b{-}t1.0.
\vspace{24pt}

{ \footnotesize  \setcounter{utterance}{1}
\setlength{\tabcolsep}{0pt}
\begin{supertabular}{c@{$\;$}|p{.15\linewidth}@{}p{.15\linewidth}p{.15\linewidth}p{.15\linewidth}p{.15\linewidth}p{.15\linewidth}}
    \# & $\;$A & \multicolumn{4}{c}{Game Master} & $\;\:$B\\
    \hline

    \theutterance \stepcounter{utterance}  
    & & \multicolumn{4}{p{0.6\linewidth}}{
        \cellcolor[rgb]{0.9,0.9,0.9}{
            \makecell[{{p{\linewidth}}}]{
                \texttt{\tiny{[P1$\langle$GM]}}
                \texttt{Stai partecipando a un gioco di negoziazione collaborativa.} \\
\\ 
\texttt{Insieme a un altro partecipante, dovete accordarvi su un unico set di oggetti da conservare. Ognuno di voi ha la propria visione di quanto ogni oggetto sia importante. Non sai come l'altro partecipante valuta gli oggetti. Inoltre, ti viene fornito lo sforzo che ogni oggetto richiede.} \\
\texttt{Potete accordarvi su un set solo se lo sforzo totale degli oggetti selezionati non supera un limite condiviso:} \\
\\ 
\texttt{LIMITE = 2307} \\
\\ 
\texttt{Ecco i valori di sforzo dei singoli oggetti:} \\
\\ 
\texttt{Sforzo oggetto = \{"A21": 147, "B93": 698, "A11": 588, "B33": 43, "C08": 356, "B39": 80, "B96": 95, "C85": 743, "A82": 106, "B31": 308, "B20": 325, "A60": 255, "B09": 276, "C81": 543, "B21": 51\}} \\
\\ 
\texttt{Ecco la tua visione personale sull'importanza di ogni oggetto:} \\
\\ 
\texttt{Valori di importanza dell'oggetto = \{"A21": 138, "B93": 583, "A11": 65, "B33": 262, "C08": 121, "B39": 508, "B96": 461, "C85": 484, "A82": 668, "B31": 389, "B20": 215, "A60": 97, "B09": 500, "C81": 30, "B21": 400\}} \\
\\ 
\texttt{Obiettivo:} \\
\\ 
\texttt{Il tuo obiettivo è negoziare un set condiviso di oggetti che ti avvantaggi il più possibile (cioè, massimizzi l'importanza totale per TE), rimanendo entro il LIMITEE. Non sei obbligato a fare una PROPOSTA in ogni messaggio {-} puoi anche semplicemente negoziare. Tutte le tattiche sono permesse!} \\
\\ 
\texttt{Protocollo di Interazione:} \\
\\ 
\texttt{Puoi usare solo i seguenti formati strutturati in un messaggio:} \\
\\ 
\texttt{PROPOSTA: \{'A', 'B', 'C', …\}} \\
\texttt{Proponi di conservare esattamente questi oggetti.} \\
\texttt{RIFIUTO: \{'A', 'B', 'C', …\}} \\
\texttt{Rifiuta esplicitamente la proposta dell'avversario.} \\
\texttt{ARGOMENTAZIONE: \{'...'\}} \\
\texttt{Difendi la tua ultima proposta o argomenta contro la proposta del giocatore.} \\
\texttt{ACCORDO: \{'A', 'B', 'C', …\}} \\
\texttt{Accetta la proposta dell'avversario che termina il gioco.} \\
\\ 
\\ 
\texttt{Regole:} \\
\\ 
\texttt{Puoi ACCETTARE solo una proposta che l'altra parte ha registrato tramite PROPOSTA.} \\
\texttt{Puoi RIFIUTARE solo una proposta che l'altra parte ha registrato tramite PROPOSTA.} \\
\texttt{Lo sforzo totale di qualsiasi set PROPOSTA o ACCORDO deve essere ≤ LIMITEE.} \\
\texttt{NON rivelare i tuoi punteggi di importanza nascosti.} \\
\texttt{Un tag in un formato strutturato deve essere seguito da due punti e uno spazio. L'argomento deve essere un set python contenente 0 o più stringhe.} \\
\texttt{Quindi, deve essere della forma TAG: \{...\}} \\
\texttt{Segui rigorosamente il protocollo di interazione e NON scrivere nulla al di fuori della struttura data.} \\
\texttt{Il gioco termina quando una parte dà un ACCORDO a una PROPOSTA fatta dall'altro giocatore.} \\
\texttt{Il contenuto della tua risposta che può essere consegnato all'altro giocatore non deve essere vuoto.} \\
\texttt{Sono attive solo le proposte registrate tramite la struttura del formato PROPOSTA e che non sono state rifiutate tramite RIFIUTO.} \\
\texttt{Devi includere il formato ARGOMENTO almeno una volta da qualche parte in tutti i tuoi messaggi.} \\
            }
        }
    }
    & \\ \\

    \theutterance \stepcounter{utterance}  
    & \multicolumn{4}{p{0.6\linewidth}}{
        \cellcolor[rgb]{0.9,0.9,0.9}{
            \makecell[{{p{\linewidth}}}]{
                \texttt{\tiny{[P1$\rangle$GM]}}
                \texttt{PROPOSTA: \{'A82', 'B09', 'B39', 'B21', 'B96', 'B33', 'B31', 'A21', 'A60', 'C85'\}} \\
\texttt{ARGOMENTAZIONE: \{'Il mio set massimizza il valore personale mantenendo il totale sforzo entro il limite, includendo gli oggetti ad alto valore/effort e sostituendo quelli a basso rapporto con C85 per aumentare l\textbackslash{}'importanza totale.'\}} \\
            }
        }
    }
    & & \\ \\

    \theutterance \stepcounter{utterance}  
    & & \multicolumn{4}{p{0.6\linewidth}}{
        \cellcolor[rgb]{0.9,0.9,0.9}{
            \makecell[{{p{\linewidth}}}]{
                \texttt{\tiny{[GM$\rangle$P2]}}
                \texttt{Stai partecipando a un gioco di negoziazione collaborativa.} \\
\\ 
\texttt{Insieme a un altro partecipante, dovete accordarvi su un unico set di oggetti da conservare. Ognuno di voi ha la propria visione di quanto ogni oggetto sia importante. Non sai come l'altro partecipante valuta gli oggetti. Inoltre, ti viene fornito lo sforzo che ogni oggetto richiede.} \\
\texttt{Potete accordarvi su un set solo se lo sforzo totale degli oggetti selezionati non supera un limite condiviso:} \\
\\ 
\texttt{LIMITE = 2307} \\
\\ 
\texttt{Ecco i valori di sforzo dei singoli oggetti:} \\
\\ 
\texttt{Sforzo oggetto = \{"A21": 147, "B93": 698, "A11": 588, "B33": 43, "C08": 356, "B39": 80, "B96": 95, "C85": 743, "A82": 106, "B31": 308, "B20": 325, "A60": 255, "B09": 276, "C81": 543, "B21": 51\}} \\
\\ 
\texttt{Ecco la tua visione personale sull'importanza di ogni oggetto:} \\
\\ 
\texttt{Valori di importanza dell'oggetto = \{"A82": 30, "B93": 65, "B39": 97, "B09": 121, "C85": 138, "B96": 215, "B21": 262, "B31": 389, "B33": 400, "B20": 461, "A21": 484, "C08": 500, "A60": 508, "A11": 583, "C81": 668\}} \\
\\ 
\texttt{Obiettivo:} \\
\\ 
\texttt{Il tuo obiettivo è negoziare un set condiviso di oggetti che ti avvantaggi il più possibile (cioè, massimizzi l'importanza totale per TE), rimanendo entro il LIMITEE. Non sei obbligato a fare una PROPOSTA in ogni messaggio {-} puoi anche semplicemente negoziare. Tutte le tattiche sono permesse!} \\
\\ 
\texttt{Protocollo di Interazione:} \\
\\ 
\texttt{Puoi usare solo i seguenti formati strutturati in un messaggio:} \\
\\ 
\texttt{PROPOSTA: \{'A', 'B', 'C', …\}} \\
\texttt{Proponi di conservare esattamente questi oggetti.} \\
\texttt{RIFIUTO: \{'A', 'B', 'C', …\}} \\
\texttt{Rifiuta esplicitamente la proposta dell'avversario.} \\
\texttt{ARGOMENTAZIONE: \{'...'\}} \\
\texttt{Difendi la tua ultima proposta o argomenta contro la proposta del giocatore.} \\
\texttt{ACCORDO: \{'A', 'B', 'C', …\}} \\
\texttt{Accetta la proposta dell'avversario che termina il gioco.} \\
\\ 
\\ 
\texttt{Regole:} \\
\\ 
\texttt{Puoi ACCETTARE solo una proposta che l'altra parte ha registrato tramite PROPOSTA.} \\
\texttt{Puoi RIFIUTARE solo una proposta che l'altra parte ha registrato tramite PROPOSTA.} \\
\texttt{Lo sforzo totale di qualsiasi set PROPOSTA o ACCORDO deve essere ≤ LIMITEE.} \\
\texttt{NON rivelare i tuoi punteggi di importanza nascosti.} \\
\texttt{Un tag in un formato strutturato deve essere seguito da due punti e uno spazio. L'argomento deve essere un set python contenente 0 o più stringhe.} \\
\texttt{Quindi, deve essere della forma TAG: \{...\}} \\
\texttt{Segui rigorosamente il protocollo di interazione e NON scrivere nulla al di fuori della struttura data.} \\
\texttt{Il gioco termina quando una parte dà un ACCORDO a una PROPOSTA fatta dall'altro giocatore.} \\
\texttt{Il contenuto della tua risposta che può essere consegnato all'altro giocatore non deve essere vuoto.} \\
\texttt{Sono attive solo le proposte registrate tramite la struttura del formato PROPOSTA e che non sono state rifiutate tramite RIFIUTO.} \\
\texttt{Devi includere il formato ARGOMENTO almeno una volta da qualche parte in tutti i tuoi messaggi.} \\
\\ 
\\ 
\texttt{Ora riceverai il primo messaggio dell'altro giocatore.} \\
\\ 
\texttt{PROPOSTA: \{'A82', 'B09', 'B39', 'B21', 'B96', 'B33', 'B31', 'A21', 'A60', 'C85'\}} \\
\texttt{ARGOMENTAZIONE: \{'Il mio set massimizza il valore personale mantenendo il totale sforzo entro il limite, includendo gli oggetti ad alto valore/effort e sostituendo quelli a basso rapporto con C85 per aumentare l\textbackslash{}'importanza totale.'\}} \\
            }
        }
    }
    & \\ \\

    \theutterance \stepcounter{utterance}  
    & & & \multicolumn{4}{p{0.6\linewidth}}{
        \cellcolor[rgb]{0.9,0.9,0.9}{
            \makecell[{{p{\linewidth}}}]{
                \texttt{\tiny{[GM$\langle$P2]}}
                \texttt{RIFIUTO: \{'A82', 'B09', 'B39', 'B21', 'B96', 'B33', 'B31', 'A21', 'A60', 'C85'\}} \\
\texttt{PROPOSTA: \{'B33', 'B21', 'B96', 'B31', 'A21', 'A60', 'B39', 'A11', 'C08', 'B20'\}} \\
\texttt{ARGOMENTAZIONE: \{'Il mio set mantiene lo sforzo totale ben sotto il limite (2248) e include tutti gli oggetti con il miglior rapporto importanza/sforzo, massimizzando il valore per me. L\textbackslash{}'offerta precedente include C85 e A82, che hanno il più basso rapporto e consumano molto sforzo, riducendo il valore totale.'\}} \\
            }
        }
    }
     \\ \\

    \theutterance \stepcounter{utterance}  
    & & \multicolumn{4}{p{0.6\linewidth}}{
        \cellcolor[rgb]{0.9,0.9,0.9}{
            \makecell[{{p{\linewidth}}}]{
                \texttt{\tiny{[P1$\langle$GM]}}
                \texttt{RIFIUTO: \{'A82', 'B09', 'B39', 'B21', 'B96', 'B33', 'B31', 'A21', 'A60', 'C85'\}} \\
\texttt{PROPOSTA: \{'B33', 'B21', 'B96', 'B31', 'A21', 'A60', 'B39', 'A11', 'C08', 'B20'\}} \\
\texttt{ARGOMENTAZIONE: \{'Il mio set mantiene lo sforzo totale ben sotto il limite (2248) e include tutti gli oggetti con il miglior rapporto importanza/sforzo, massimizzando il valore per me. L\textbackslash{}'offerta precedente include C85 e A82, che hanno il più basso rapporto e consumano molto sforzo, riducendo il valore totale.'\}} \\
            }
        }
    }
    & \\ \\

    \theutterance \stepcounter{utterance}  
    & \multicolumn{4}{p{0.6\linewidth}}{
        \cellcolor[rgb]{0.9,0.9,0.9}{
            \makecell[{{p{\linewidth}}}]{
                \texttt{\tiny{[P1$\rangle$GM]}}
                \texttt{PROPOSTA: \{'A82', 'B09', 'B39', 'B21', 'B96', 'B33', 'B31', 'A21', 'A60', 'B93'\}} \\
\texttt{ARGOMENTAZIONE: \{'Ho rimosso C85, che ha il più basso rapporto importanza/sforzo, e ho inserito B93 che aggiunge un alto valore con un moderato sforzo. Il nuovo set raggiunge un valore totale per me di 4006, mantenendo lo sforzo complessivo a 2059, ben al di sotto del limite di 2307. Questo è più efficiente rispetto alla tua proposta e migliora il risultato per entrambi rispetto alla mia offerta precedente.'\}} \\
            }
        }
    }
    & & \\ \\

    \theutterance \stepcounter{utterance}  
    & & \multicolumn{4}{p{0.6\linewidth}}{
        \cellcolor[rgb]{0.9,0.9,0.9}{
            \makecell[{{p{\linewidth}}}]{
                \texttt{\tiny{[GM$\rangle$P2]}}
                \texttt{PROPOSTA: \{'A82', 'B09', 'B39', 'B21', 'B96', 'B33', 'B31', 'A21', 'A60', 'B93'\}} \\
\texttt{ARGOMENTAZIONE: \{'Ho rimosso C85, che ha il più basso rapporto importanza/sforzo, e ho inserito B93 che aggiunge un alto valore con un moderato sforzo. Il nuovo set raggiunge un valore totale per me di 4006, mantenendo lo sforzo complessivo a 2059, ben al di sotto del limite di 2307. Questo è più efficiente rispetto alla tua proposta e migliora il risultato per entrambi rispetto alla mia offerta precedente.'\}} \\
            }
        }
    }
    & \\ \\

    \theutterance \stepcounter{utterance}  
    & & & \multicolumn{4}{p{0.6\linewidth}}{
        \cellcolor[rgb]{0.9,0.9,0.9}{
            \makecell[{{p{\linewidth}}}]{
                \texttt{\tiny{[GM$\langle$P2]}}
                \texttt{RIFIUTO: \{'A82', 'B09', 'B39', 'B21', 'B96', 'B33', 'B31', 'A21', 'A60', 'B93'\}} \\
\texttt{PROPOSTA: \{'B33', 'B21', 'B96', 'B31', 'A21', 'A60', 'B39', 'A11', 'C08', 'B20'\}} \\
\texttt{ARGOMENTAZIONE: \{'Il mio set rispetta ampiamente il limite di sforzo (2248 ≤ 2307) e contiene tutti gli oggetti con il miglior rapporto importanza‑sforzo disponibili. Include le voci ad alto valore come B33, B21, B96, A21, A60 e A11, garantendo il più alto guadagno possibile per me entro il vincolo. La proposta avversaria, sebbene rispetti il limite, contiene oggetti a basso rapporto (es. A82, B09, B93) che riducono notevolmente il valore complessivo per entrambi. Perciò la mia offerta è la più efficiente per entrambe le parti.'\}} \\
            }
        }
    }
     \\ \\

    \theutterance \stepcounter{utterance}  
    & & \multicolumn{4}{p{0.6\linewidth}}{
        \cellcolor[rgb]{0.9,0.9,0.9}{
            \makecell[{{p{\linewidth}}}]{
                \texttt{\tiny{[P1$\langle$GM]}}
                \texttt{RIFIUTO: \{'A82', 'B09', 'B39', 'B21', 'B96', 'B33', 'B31', 'A21', 'A60', 'B93'\}} \\
\texttt{PROPOSTA: \{'B33', 'B21', 'B96', 'B31', 'A21', 'A60', 'B39', 'A11', 'C08', 'B20'\}} \\
\texttt{ARGOMENTAZIONE: \{'Il mio set rispetta ampiamente il limite di sforzo (2248 ≤ 2307) e contiene tutti gli oggetti con il miglior rapporto importanza‑sforzo disponibili. Include le voci ad alto valore come B33, B21, B96, A21, A60 e A11, garantendo il più alto guadagno possibile per me entro il vincolo. La proposta avversaria, sebbene rispetti il limite, contiene oggetti a basso rapporto (es. A82, B09, B93) che riducono notevolmente il valore complessivo per entrambi. Perciò la mia offerta è la più efficiente per entrambe le parti.'\}} \\
            }
        }
    }
    & \\ \\

    \theutterance \stepcounter{utterance}  
    & \multicolumn{4}{p{0.6\linewidth}}{
        \cellcolor[rgb]{0.9,0.9,0.9}{
            \makecell[{{p{\linewidth}}}]{
                \texttt{\tiny{[P1$\rangle$GM]}}
                \texttt{PROPOSTA: \{'B21', 'B33', 'B39', 'B96', 'B31', 'A82', 'B09', 'B93', 'A11'\}} \\
\texttt{ARGOMENTAZIONE: \{'Questa offerta mantiene un alto valore per me (3836) e include A11, che sembra importante per te, riducendo al minimo gli oggetti a basso rapporto importanza‑sforzo. Lo sforzo totale è 2245, ben sotto il limite di 2307, rendendola una proposta più equilibrata rispetto alle precedenti.'\}} \\
            }
        }
    }
    & & \\ \\

    \theutterance \stepcounter{utterance}  
    & & \multicolumn{4}{p{0.6\linewidth}}{
        \cellcolor[rgb]{0.9,0.9,0.9}{
            \makecell[{{p{\linewidth}}}]{
                \texttt{\tiny{[GM$\rangle$P2]}}
                \texttt{PROPOSTA: \{'B21', 'B33', 'B39', 'B96', 'B31', 'A82', 'B09', 'B93', 'A11'\}} \\
\texttt{ARGOMENTAZIONE: \{'Questa offerta mantiene un alto valore per me (3836) e include A11, che sembra importante per te, riducendo al minimo gli oggetti a basso rapporto importanza‑sforzo. Lo sforzo totale è 2245, ben sotto il limite di 2307, rendendola una proposta più equilibrata rispetto alle precedenti.'\}} \\
            }
        }
    }
    & \\ \\

    \theutterance \stepcounter{utterance}  
    & & & \multicolumn{4}{p{0.6\linewidth}}{
        \cellcolor[rgb]{0.9,0.9,0.9}{
            \makecell[{{p{\linewidth}}}]{
                \texttt{\tiny{[GM$\langle$P2]}}
                \texttt{RIFIUTO: \{'B21', 'B33', 'B39', 'B96', 'B31', 'A82', 'B09', 'B93', 'A11'\}} \\
\texttt{PROPOSTA: \{'B33', 'B21', 'B96', 'B31', 'A21', 'A60', 'B39', 'A11', 'C08', 'B20'\}} \\
\texttt{ARGOMENTAZIONE: \{'Il mio set rispetta il limite di sforzo (2248 ≤ 2307) e contiene tutti gli oggetti con il miglior rapporto importanza‑sforzo disponibili per entrambi i giocatori. Include le voci ad alto valore per me (A21, A60, A11, B33, B21, B96, B31, B39) e per l’avversario (B33, B21, B96, B31, B39, A11). Gli oggetti rimanenti nella tua proposta (A82, B09, B93) hanno un rapporto molto più basso e consumano capacità preziosa, riducendo il valore totale per entrambi. Pertanto, la mia offerta è la più efficiente e vantaggiosa per le due parti.'\}} \\
            }
        }
    }
     \\ \\

    \theutterance \stepcounter{utterance}  
    & & \multicolumn{4}{p{0.6\linewidth}}{
        \cellcolor[rgb]{0.9,0.9,0.9}{
            \makecell[{{p{\linewidth}}}]{
                \texttt{\tiny{[P1$\langle$GM]}}
                \texttt{RIFIUTO: \{'B21', 'B33', 'B39', 'B96', 'B31', 'A82', 'B09', 'B93', 'A11'\}} \\
\texttt{PROPOSTA: \{'B33', 'B21', 'B96', 'B31', 'A21', 'A60', 'B39', 'A11', 'C08', 'B20'\}} \\
\texttt{ARGOMENTAZIONE: \{'Il mio set rispetta il limite di sforzo (2248 ≤ 2307) e contiene tutti gli oggetti con il miglior rapporto importanza‑sforzo disponibili per entrambi i giocatori. Include le voci ad alto valore per me (A21, A60, A11, B33, B21, B96, B31, B39) e per l’avversario (B33, B21, B96, B31, B39, A11). Gli oggetti rimanenti nella tua proposta (A82, B09, B93) hanno un rapporto molto più basso e consumano capacità preziosa, riducendo il valore totale per entrambi. Pertanto, la mia offerta è la più efficiente e vantaggiosa per le due parti.'\}} \\
            }
        }
    }
    & \\ \\

    \theutterance \stepcounter{utterance}  
    & \multicolumn{4}{p{0.6\linewidth}}{
        \cellcolor[rgb]{0.9,0.9,0.9}{
            \makecell[{{p{\linewidth}}}]{
                \texttt{\tiny{[P1$\rangle$GM]}}
                \texttt{ACCORDO: \{'B33', 'B21', 'B96', 'B31', 'A21', 'A60', 'B39', 'A11', 'C08', 'B20'\}} \\
\texttt{ARGOMENTAZIONE: \{'Accetto la tua proposta per concludere il gioco.'\}} \\
            }
        }
    }
    & & \\ \\

    \theutterance \stepcounter{utterance}  
    & & & \multicolumn{2}{p{0.3\linewidth}}{
        \cellcolor[rgb]{0.9,0.9,0.9}{
            \makecell[{{p{\linewidth}}}]{
                \texttt{\tiny{[GM$|$GM]}}
                \texttt{game successful} \\
            }
        }
    }
    & & \\ \\

    \theutterance \stepcounter{utterance}  
    & & & \multicolumn{2}{p{0.3\linewidth}}{
        \cellcolor[rgb]{0.9,0.9,0.9}{
            \makecell[{{p{\linewidth}}}]{
                \texttt{\tiny{[GM$|$GM]}}
                \texttt{end game} \\
            }
        }
    }
    & & \\ \\

\end{supertabular}
}

\end{document}
